% MODERNIZATION_CASE_STUDY.tex — XBoing: A Case Study in Agentic Software Modernization
%
% Intended audience: engineers evaluating agentic software-engineering
% methods for modernizing legacy code, with emphasis on the punt-labs
% tooling used here (ethos mission contracts, z-spec formal methods).
%
% Figures come from docs/metrics.tex (\Metric* macros), generated from
% docs/metrics.json — never hand-edit numbers here; edit the JSON. The ADR
% table comes from docs/adr_table.tex, generated from docs/DESIGN.md.
% Build everything (regenerate the derived tables + compile all report PDFs)
% with `make docs-pdf`. Commit the .tex, the generated tables, and the PDFs
% together.

\documentclass[11pt,a4paper]{article}

% Shared preamble (packages, macros, listings style) — see docs/xboing-doc.sty
\usepackage{xboing-doc}

% Project figures, generated from docs/metrics.json by
% scripts/gen_metrics_tex.sh. Never hand-edit metrics here — edit the JSON
% and regenerate so every figure stays in sync with its source.
% GENERATED by scripts/gen_metrics_tex.sh from docs/metrics.json — do not edit by hand.
\newcommand{\MetricVersion}{1.0.11}
\newcommand{\MetricFirstCommit}{2025-11-26}
\newcommand{\MetricLastCommit}{2026-09-26}
\newcommand{\MetricActiveDays}{54}
\newcommand{\MetricMergedPRs}{208}
\newcommand{\MetricCommits}{354}
\newcommand{\MetricADRs}{74}
\newcommand{\MetricSrcLines}{19{,}820}
\newcommand{\MetricSrcFiles}{47}
\newcommand{\MetricHdrLines}{2{,}993}
\newcommand{\MetricHdrFiles}{51}
\newcommand{\MetricTestLines}{21{,}742}
\newcommand{\MetricTestFiles}{66}
\newcommand{\MetricLegacyLines}{15{,}405}
\newcommand{\MetricLegacyFiles}{32}
\newcommand{\MetricCmockaCases}{1{,}337}
\newcommand{\MetricCtestSuites}{51}
\newcommand{\MetricBeadsClosed}{240}
\newcommand{\MetricBeadsTotal}{261}
\newcommand{\MetricLevels}{80}
\newcommand{\MetricClosedMissions}{107}


% --- Title -------------------------------------------------------------------
\title{%
    \textbf{XBoing: A Case Study in Agentic Software Modernization}\\[0.5em]
    \large Reviving a 1996 X11 Game with LLM Agents, Mission Contracts,\\
    \large and Formal Methods
}
\author{XBoing Modernization Project}
\date{September 2026}

\begin{document}
\maketitle
\sloppy % path-heavy technical prose; prefer loose spacing over overfull boxes

\begin{abstract}
\noindent
XBoing is a fast, colorful breakout arcade game Justin C.\ Kibell wrote in
pure Xlib between 1993 and 1996. This document is a case study of bringing
it forward to a modern SDL2 codebase --- not by rewriting it in the loose
sense, but by reimplementing it from scratch under a hard constraint:
100\% behavioral fidelity to the 1996 original, with every deviation
limited and documented. The distinguishing feature of the project is not
the port itself but the \emph{method}: the work was driven by LLM agents
operating under a principal-engineer standard, coordinated through
\emph{ethos} mission contracts (typed delegation with a worker and a
distinct evaluator), and hardened in places with \emph{z-spec} formal
methods (Z specifications model-checked with \code{probcli}). This report
documents what was built, how the agentic method worked in practice, the
two punt-labs tools that made it auditable and verifiable, and --- in
keeping with the project's ``evidence, not nostalgia'' mandate --- where
it was hard and what recurred.
\end{abstract}

\tableofcontents
\newpage

% =============================================================================
\section{Introduction}
% =============================================================================

\subsection{What was built}

XBoing 2.4 shipped on 22 November 1996 as pure Xlib --- no Motif, no Xt ---
with XPM pixmap graphics, \code{fork()}-and-\code{/dev/dsp} audio across
twelve platform-specific drivers, and compile-time paths baked in with
\code{\#define}. It then sat unmaintained for roughly two decades.

The modernization produces a current SDL2 game (release 1.0.11) that
installs via Homebrew and Debian packaging and plays on contemporary Linux
and macOS. Two source trees coexist in the repository and share
\emph{zero} lines:

\begin{itemize}[nosep]
  \item \filename{original/} --- Kibell's 1996 sources, preserved verbatim
        as the canonical reference. Never edited; cited line-by-line
        (\filename{original/<file>.c:<line>}) whenever behavior is
        reproduced.
  \item \filename{src/} + \filename{include/} --- the new SDL2
        implementation. The shipping game.
\end{itemize}

\subsection{Reimplementation, not reinterpretation}

The precise claim is worth stating carefully, because two common framings
are both wrong. It is not a line-level port (nothing in \filename{src/} is
transliterated from \filename{original/}), and it is not a rewrite in the
sense of reimagining the game. It is a faithful \emph{reimplementation from
scratch}: new code, same game. The 1996 behavior is treated as a
specification to reproduce exactly --- level file format, scoring values,
physics constants (\code{MAX\_BALLS=5}, the 18$\times$9 grid, paddle bounce
trigonometry, \code{DIST\_BASE=30}), timing feel, the sixteen-mode state
machine. All 80 original levels load and play unchanged.

Fidelity is the default; deviations are limited and documented as
Architecture Decision Records (ADRs). Two examples: save/load v2 resumes a
ball mid-flight rather than resetting it to the paddle
(\adref{038}); and \code{PADDLE\_VELOCITY} is a known, tracked deviation
(paddle moves faster than the original) recorded in the gameplay-parity
audit. ``Faithful by default, deviations only when forced and always
documented'' is the operative rule, not a flat claim of perfection.

\subsection{The architecture shift}

The port replaces every platform boundary while preserving the game
logic. \Cref{tab:fromto} summarizes the substitutions; the companion
documents \filename{docs/ARCHITECTURE\_LEGACY.tex} and
\filename{docs/ARCHITECTURE\_MODERN.tex} give the full architectural
treatment.

\begin{table}[h]
\centering
\caption{Platform substitutions (game logic preserved)}
\label{tab:fromto}
\begin{tabularx}{\textwidth}{llX}
\toprule
\textbf{Subsystem} & \textbf{Original (1996)} & \textbf{Modernized} \\
\midrule
Windowing   & Xlib + private colormap + GCs        & SDL2 \\
Graphics    & XPM via libXpm                        & PNG via SDL2\_image \\
Audio       & \code{fork()} + \code{/dev/dsp} + 12 drivers & SDL2\_mixer \\
Sound format& \code{.au} (Sun audio)               & OGG / WAV \\
Build       & hand-written Makefile + Imake         & CMake \\
Paths       & compile-time \code{\#define}s          & XDG Base Directory spec \\
Fonts       & XLFD bitmap fonts                     & bundled TTF via SDL2\_ttf \\
Collision   & \code{XPolygonRegion} (server-side)   & pure-C point-in-triangle \\
\bottomrule
\end{tabularx}
\end{table}

% =============================================================================
\section{By the numbers}
% =============================================================================

All figures are sourced from \filename{docs/metrics.json}, regenerated from
\code{git log}, \code{bd stats}, and \code{cloc}. Line counts are reported
as code lines (\code{cloc}, excluding comments and blanks); raw
\code{wc -l} totals are larger.

\begin{table}[h]
\centering
\caption{Project metrics at release \MetricVersion{} (snapshot: \code{origin/master})}
\label{tab:metrics}
\begin{tabularx}{\textwidth}{Xr}
\toprule
\textbf{Metric} & \textbf{Value} \\
\midrule
Active development                 & \MetricActiveDays{} days (part-time, \MetricFirstCommit{} → \MetricLastCommit{}) \\
Merged pull requests              & \MetricMergedPRs{} \\
Commits on \code{master}          & \MetricCommits{} \\
Architecture Decision Records     & \MetricADRs{} \\
Modern source (\filename{src/})   & \MetricSrcLines{} code lines / \MetricSrcFiles{} files \\
Modern headers (\filename{include/}) & \MetricHdrLines{} code lines / \MetricHdrFiles{} files \\
Test code (\filename{tests/})     & \MetricTestLines{} code lines / \MetricTestFiles{} files \\
Legacy reference (\filename{original/}) & \MetricLegacyLines{} code lines / \MetricLegacyFiles{} \code{.c} files \\
\code{cmocka} test cases          & \MetricCmockaCases{} \\
CTest suites                      & \MetricCtestSuites{} \\
Issues (beads) closed / total     & \MetricBeadsClosed{} / \MetricBeadsTotal{} \\
Playable levels                   & \MetricLevels{} \\
\bottomrule
\end{tabularx}
\end{table}

Two ratios are worth calling out. The test suite is larger than the code it
exercises: \MetricTestLines{} code lines of tests against \MetricSrcLines{} of modern source,
a ratio of $\approx$1.1:1 (1.2:1 by raw line count). And the modern
codebase is larger than the 1996 original it reproduces --- the safety
apparatus (opaque contexts, injected callbacks, exhaustive tests) costs
lines, and buys the ability to change the game without fear.

For contrast, the first ``modernization complete'' milestone was recorded
on 26 February 2026 at 64 merged PRs, 91 closed issues, and 940 tests ---
roughly one quarter of the eventual total. The bulk of the fidelity work
(the parity audits, the formal-methods hardening, the agentic tooling
adoption) happened \emph{after} the code first ``worked.''

% =============================================================================
\section{Method: agentic software engineering}
% =============================================================================

\subsection{The standard}

The governing instruction is that every change leaves the codebase better
than it was found, with root causes proven by facts and tests rather than
asserted. Crucially, the human operator does not write the production code
and does not let a single agent both write and bless its own work.
Delegation is structured, typed, and auditable --- which is what the
\emph{ethos} tooling provides.

\subsection{Roles}

Work is organized around three roles:

\begin{description}[nosep]
  \item[Leader] (\texttt{claude}) writes the contract and reviews. It does
        not write production code solo. For non-code work (design, reports,
        reviews, investigations) it may act as the worker.
  \item[Worker] a domain specialist that executes the contract's write-set.
  \item[Evaluator] a \emph{distinct} agent ($\neq$ worker, $\neq$ leader)
        that reviews the result adversarially.
\end{description}

The domain specialists are modeled on the people who would actually own
each concern (\cref{tab:personas}).

\begin{table}[h]
\centering
\caption{Expert personas}
\label{tab:personas}
\begin{tabularx}{\textwidth}{llX}
\toprule
\textbf{Handle} & \textbf{Persona} & \textbf{Domain} \\
\midrule
\texttt{jck} & Justin C.\ Kibell   & Original author; game feel, physics, scoring. Approves gameplay-affecting change. \\
\texttt{jdc} & John D.\ Carmack     & Modern C, sanitizers; primary C implementer. \\
\texttt{sjl} & Sam J.\ Lantinga     & SDL2 author; rendering and audio port. \\
\texttt{gjm} & Glenford J.\ Myers   & Testing; harness design, extracting pure functions. \\
\bottomrule
\end{tabularx}
\end{table}

\subsection{Scale}

Across the project, \MetricClosedMissions{} closed missions are recorded in the tracked ledger
\filename{.punt-labs/ethos/missions.jsonl}. By archetype: 100
\emph{implement}, 6 \emph{design}, 2 \emph{investigate}, 1 \emph{report}.
The worker$\to$evaluator pairings that actually ran show the specialist
model in practice: \texttt{jdc}$\to$\texttt{jck} (39),
\texttt{gjm}$\to$\texttt{jdc} (27), \texttt{jdc}$\to$\texttt{gjm} (23),
\texttt{sjl}$\to$\texttt{jdc} (9), and smaller counts. C implementation
reviewed by the tester; the tester's harnesses reviewed by the C engineer;
anything that touches game feel routed to the original author for approval.

% =============================================================================
\section{Tooling: ethos (mission contracts)}
% =============================================================================

\subsection{A mission is a contract, not a chat}

The load-bearing idea is that delegated work runs as a typed contract, not
an ad-hoc agent spawn. A contract pins the worker, the evaluator (with a
content hash, so the evaluator cannot be silently swapped), the
\emph{write-set} of files the worker may touch, and machine-checkable
success criteria. A pre-tool hook blocks edits outside the write-set at
runtime, so scope creep fails closed.

The protocol is \code{create} $\to$ \code{result} $\to$ \code{reflect}
$\to$ \code{advance} $\to$ \code{close}. The store refuses self-review,
refuses closing without a submitted result, and refuses advancing a round
without a recorded reflection --- the process cannot be shortcut.

\subsection{A real contract}

\Cref{lst:contract} is an actual committed contract:
\missionid{m-2026-07-12-016}, the ammo-belt fix. Note that its success
criteria are concrete and testable, and that one of them is literally a
formal invariant --- this is the seam where ethos and z-spec meet.

\begin{lstlisting}[caption={The ammo-belt fix contract (excerpt, verbatim)}, label={lst:contract}]
mission_id: m-2026-07-12-016
type: implement
leader: claude
worker: sjl
evaluator:
    handle: jdc
    pinned_at: "2026-07-12T16:05:18Z"
    hash: c08223356d3e6d5aed56e28da477a1aa2707509febfef7c5f62c25871a3e841f
write_set:
    - src/game_render.c
    - include/game_render.h
success_criteria:
    - Remove BOTH the unlimited early-return AND the GUN_MAX_AMMO cap from
      game_render_ammo_belt -- removing only one still desyncs (jck).
    - Extract a pure count seam game_render_ammo_belt_count(ctx) that the
      test pins (belt count == bonus bullet_count).
    - Cite original/level.c:324-334 and original/blocks.c:1588-1593.
    - cmake build zero-warning; ctest stays green.
budget: { rounds: 3, reflection_after_each: true }
context: |
    Bug xboing-65h -- ammo-belt desync. The belt hid when unlimited and
    capped at GUN_MAX_AMMO, desyncing from the end-of-level bonus which
    reads raw ammo. jck approved the faithful fix. Formal model
    docs/bullet_ammo.tex pins the invariant displayed == actual == bonus.
\end{lstlisting}

\subsection{The reflect/advance loop does real work}

The evaluator is adversarial, and it sends work back. Mission
\missionid{m-2026-07-11-011} records the full lifecycle: the evaluator
returned \textsc{changes-requested} with three findings, forcing a second
round before the result passed. The event log is append-only:

\begin{lstlisting}[caption={Mission lifecycle events (log.jsonl)}, label={lst:lifecycle}]
{"event":"create","actor":"claude"}
{"event":"reflect","details":{"converging":false,"recommendation":"continue","round":1}}
{"event":"round_advanced","details":{"from_round":1,"to_round":2}}
{"event":"result","details":{"verdict":"pass","round":2}}
{"event":"close","details":{"round":2,"status":"closed","verdict":"pass"}}
\end{lstlisting}

\subsection{The audit chain}

Closing a mission appends a trace line to the git-tracked ledger, and
commits carry \code{Mission:} and \code{Delegation:} trailers. The intended
reconstruction is \code{git blame} $\to$ commit $\to$ trailer $\to$
contract $\to$ prompt $\to$ full audit trail: every shipped line traces
back to the contract that authorized it and the review that blessed it. On
\code{master} there are 45 \code{Mission:} trailers and 11
\code{Delegation:} trailers.

\subsection{An honest caveat}

The audit mechanism was itself dogfooded, and its failure was caught rather
than hidden. When ethos was first enabled, the commit-message hook that
should have auto-stamped the trailers stamped nothing --- a
session-binding bug (tracked upstream) meant commits made under an open
mission carried no trailer. The trailers now present were added manually
after the bug was diagnosed and documented (PR \#220, mission
\missionid{m-2026-09-06-009}). This is exactly the kind of failure the
case study exists to record: the tool that guarantees traceability was
found untrustworthy, empirically, by using it.

% =============================================================================
\section{Tooling: z-spec (formal methods)}
% =============================================================================

\subsection{Why formal methods on a breakout game}

Most of XBoing is not a candidate for formal specification --- geometry and
rendering are better served by characterization tests. But a handful of
properties are (a) invariants the player can observe when they break, and
(b) reachable only through specific state sequences that are easy to miss
in hand-written tests. For those, the project used z-spec: a Z
specification, type-checked with \code{fuzz} and model-checked with
\code{probcli}, with each safety property proven by model-checking toward
its \emph{negation}. A reachable state matching the negation is a
counter-example that proves the defect; ``no counter-example, all states
visited'' after the fix proves the guard closes it. The technique was
applied selectively --- several ADRs explicitly record \emph{not} using it
where the logic is deterministic and stateless.

\subsection{The ammo-belt desync}

The specification \filename{docs/bullet\_ammo.tex} models three quantities
the code keeps distinct --- the authoritative ammo count, the number the
HUD belt draws, and the end-of-level bonus --- and pins the invariant
\code{displayedAmmo = actualAmmo}. Model-checking the \emph{as-coded} belt
rule, the property that the bonus and the display agree excluded 441 of 486
reachable states: the entire desync region. The counter-example was
concrete --- pick up a max-ammo block, enter the bonus, and the belt showed
zero while the bonus paid twenty-one. Under the fix (draw the raw count, no
unlimited special-case, no cap), the model covers 45 of 45 states with the
invariant preserved by every operation. The fix is recorded in \adref{063},
and the same invariant is pinned in the code by a red$\to$green regression
test. This is the property named in the contract of \cref{lst:contract}:
the formal model, the contract's success criterion, and the regression
test are the same statement in three forms.

\subsection{The screen state machine}

The specification \filename{docs/specs/2026-07-04-screen-state-machine.tex}
models the flat mode machine and its shared resources, with three safety
invariants (stated deliberately as observer schemas so that violating
transitions stay reachable rather than being silently disabled):

\begin{lstlisting}[caption={Screen-machine safety invariants (Z)}, label={lst:zschema}]
SafeGame:      mode = mGame => blocksLoaded = true
SafeHighscore: (mode = mHighscore and boardChoice = bGlobal
                 and globalHasData = false) => personalHasData = false
SafeAttract:   (mode in attractModes and mode /= mHighscore)
                 => gameActive = false
\end{lstlisting}

Each invariant caught a real defect, proven reachable in the as-coded model
and proven unreachable after the fix in one exhaustive pass:
the empty Hall-of-Fame on unprivileged installs (\adref{054},
\code{SafeHighscore}, 138 states after fix); a \code{game\_active} flag
leaking out of the attract cycle (\adref{055}, \code{SafeAttract}, a
7-step counter-example trace, 90 states after fix); and pressing Space
landing in the bonus screen with an empty grid after a silent level-load
failure (\adref{056}, \code{SafeGame}, 76 states after fix). Each fixed
invariant is both a re-proven model postcondition and a named
red$\to$green regression test.

\subsection{A caveat on the artifacts}

One committed artifact,
\filename{docs/specs/2026-07-04-screen-state-machine.report.json}, is a raw
\code{zspec test} run at a small set size that deadlocks --- it is a
tooling snapshot, distinct from the goal-directed proofs above. The
authoritative results are the ``all states visited'' outcomes in the
specification's verification section and the ADRs, not that JSON.

% =============================================================================
\section{Testing and fidelity}
% =============================================================================

The test strategy is a five-layer pyramid: (1) \code{cmocka} unit tests of
pure logic; (2) integration tests running the game loop headless with the
SDL dummy drivers; (3) replay/golden tests of full input sequences; (4)
visual-fidelity tests that capture screenshots and compare them against the
1996 binary, including an LLM-judged comparison; and (5) parser fuzzing.
The savegame-v2 fixture mechanism (\adref{038}) is load-bearing here: it
lets a test seed arbitrary deep game state (level 25, a specific block
mid-cooldown, an in-flight ball) directly, instead of replaying inputs to
reach it.

Fidelity was verified subsystem by subsystem through six parity audits ---
audio volume, editor, gameplay logic, animation cadence, key commands, and
sound effects --- each comparing the modern build against
\filename{original/} and filing discrepancies as beads. Not every audit
finding was a real bug: at least one (bonus-screen tally ``too slow'') was
a false positive, and the pacing turned out to be an intentional
readability trade-off. Recording the false positives is part of the
evidence.

% =============================================================================
\section{What was hard, and what recurred}
% =============================================================================

In keeping with the project's mandate --- evidence, not nostalgia --- the
honest retrospective matters as much as the wins.

\begin{itemize}[nosep]
  \item \textbf{Documentation drift.} The retrospective docs froze at the
        February ``100\% complete'' milestone and did not track the large
        body of work that landed after it; this very report is the
        correction. The lesson: a single machine-readable metrics source
        (\filename{docs/metrics.json}) beats prose numbers that silently
        rot.
  \item \textbf{The traceability tool's own bug.} The trailer-stamping
        hook was broken on first enable (\S5.6). Dogfitting the audit
        mechanism is what surfaced it.
  \item \textbf{Config contradictions.} Three project files independently
        encoded the same facts and drifted apart, requiring a
        reconciliation pass (PR \#222). Single-source-of-truth is a
        recurring failure mode, not a one-off.
  \item \textbf{Faithful-looking bugs that are not bugs.} Several
        ``deviations'' (the ammo morph-block scoring zero, the
        $\approx$22.5s auto-launch delay, the bonus tally pacing) were
        verified \emph{faithful} to 1996 and deliberately kept. Fidelity
        work is as much about resisting spurious fixes as making real ones.
\end{itemize}

% =============================================================================
\section{Reproducibility}
% =============================================================================

Every claim here resolves to a committed artifact:

\begin{itemize}[nosep]
  \item Metrics: \filename{docs/metrics.json} (with the regenerating command per field).
  \item Mission contracts and the audit ledger:
        \filename{.punt-labs/ethos/missions/} and
        \filename{.punt-labs/ethos/missions.jsonl}.
  \item Formal specifications: \filename{docs/bullet\_ammo.tex} and
        \filename{docs/specs/2026-07-04-screen-state-machine.tex}.
  \item Decisions: \filename{docs/DESIGN.md} (\adref{001} through \adref{074}).
  \item Architecture: \filename{docs/ARCHITECTURE\_LEGACY.tex} and
        \filename{docs/ARCHITECTURE\_MODERN.tex}.
\end{itemize}

\vspace{1em}
\noindent\textbf{Acknowledgement.} XBoing was created by Justin C.\ Kibell
(1993--1997). The 1996 sources are preserved verbatim in
\filename{original/} under their original permissive license.

\end{document}
